% Конспект-методичка лекции 10. Сборка из корня:
% latexmk 10/ascent-to-orbit.tex
\documentclass[a4paper,11pt]{article}
\usepackage{iftex}
\RequirePDFTeX
\usepackage{cmap}
\usepackage[T2A]{fontenc}
\usepackage[utf8]{inputenc}
\usepackage[english,russian]{babel}
\usepackage{PTSerif}
\usepackage{amsmath,amssymb,amsthm,mathtools}
\usepackage[a4paper,left=20mm,right=20mm,top=22mm,bottom=23mm,headheight=14pt]{geometry}
\usepackage{microtype,anyfontsize}
\usepackage{graphicx,xcolor,tikz,booktabs,tabularx,array,longtable}
\usetikzlibrary{arrows.meta,calc,positioning}
\usepackage{enumitem,fancyhdr,needspace,float,placeins}
\usepackage[most]{tcolorbox}
\definecolor{accent}{HTML}{187AA5}
\definecolor{muted}{HTML}{5A6166}
\definecolor{figink}{HTML}{25323F}
\definecolor{figtan}{HTML}{087F78}
\definecolor{figorange}{HTML}{B57714}
\definecolor{figpurple}{HTML}{79539C}
\newtcolorbox{explanation}[1]{enhanced,
 colback=accent!3!white,colframe=accent!80!black,
 boxrule=.5pt,toprule=1.2pt,arc=1mm,
 left=3mm,right=3mm,top=3mm,bottom=2mm,
 before skip=9pt,after skip=9pt,
 title={#1},fonttitle=\bfseries,coltitle=accent!70!black,
 attach boxed title to top left={xshift=3mm,yshift=-2mm},
 boxed title style={colback=white,frame hidden,boxrule=0pt,
 left=1.5mm,right=1.5mm,top=.4mm,bottom=.4mm}}
\newtcolorbox{remarkbox}{enhanced,
 colback=orange!4!white,colframe=orange!65!black,
 boxrule=.5pt,toprule=1.2pt,arc=1mm,
 left=3mm,right=3mm,top=2mm,bottom=2mm,
 before skip=9pt,after skip=9pt}
\usepackage{hyperref}
\usepackage{bookmark}
\hypersetup{unicode=true,colorlinks=true,linkcolor=accent,citecolor=accent,urlcolor=accent,
 pdftitle={Задача выведения КА на орбиту},
 pdfauthor={Материалы курса «Введение в маневрирование КА»},
 pdfsubject={Динамика ракеты, гравитационный разворот, потери, довыведение и пристрелка}}
\newcommand{\dd}{\mathop{}\!\mathrm d}
\newcommand{\vect}[1]{\boldsymbol{#1}}
\newcommand{\norm}[1]{\left\lVert#1\right\rVert}
\DeclareMathOperator{\atanTwo}{atan2}
\numberwithin{equation}{section}
\setcounter{tocdepth}{2}
\setcounter{secnumdepth}{2}
\setlength{\parindent}{1.3em}
\setlength{\parskip}{2pt}
\linespread{1.055}
\setlist{itemsep=2pt,topsep=4pt,parsep=0pt}
\setlength{\emergencystretch}{1.5em}
\clubpenalty=10000
\widowpenalty=10000
\brokenpenalty=10000
\displaywidowpenalty=10000
\allowdisplaybreaks[1]
\setcounter{topnumber}{2}
\renewcommand{\topfraction}{.88}
\renewcommand{\textfraction}{.10}
\renewcommand{\floatpagefraction}{.82}
\makeatletter
\setlength{\@fptop}{0pt}
\makeatother
\pagestyle{fancy}
\fancyhf{}
\fancyhead[L]{\small\color{muted}Задача выведения КА на орбиту}
\fancyfoot[C]{\small\thepage}
\renewcommand{\headrulewidth}{.3pt}
\newcommand{\coursefigure}[3][.98]{%
\begin{figure}[!htbp]\centering
\includegraphics[width=#1\textwidth]{images/#2.pdf}
\caption{#3}\label{fig:#2}\end{figure}}
\newcommand{\courseplot}[3][.98]{%
\begin{figure}[H]\centering
\includegraphics[width=#1\textwidth]{images/#2.pdf}
\caption{#3}\label{fig:#2}\end{figure}}
\begin{document}
\begin{titlepage}
\centering
{\large Московский физико-технический институт\par}
\vspace{8mm}
{\large Курс «Введение в маневрирование КА»\par}
\vspace{30mm}
{\color{accent}\rule{28mm}{.8pt}\par}
\vspace{9mm}
{\LARGE\bfseries Задача выведения КА\\[3mm]на орбиту\par}
\vspace{11mm}
{\large Методические материалы к лекции 10\par}
\vspace{8mm}
{Конечное орбитальное состояние. Динамика ракеты.\\
Программа выведения, ограничения и довыведение.\\
Расчёт траектории и уточнение управления.\par}
\vfill
{Рабочая редакция\par}
\vspace{3mm}
{14 ноября 2026 года\par}
\end{titlepage}
\setcounter{page}{2}
\begingroup\small
\setlength{\parskip}{0pt}
\tableofcontents
\endgroup
\clearpage

\section*{Предмет рассмотрения и обозначения}
\addcontentsline{toc}{section}{Предмет рассмотрения и обозначения}

Задача выведения --- сформировать положение и инерциальную скорость
КА, соответствующие заданной орбите. Рассмотрим, как эту цель
связывают с тягой, расходом топлива и программой движения ракеты.
Основу изложения составляют разделы Teofilatto об атмосферном
и внеатмосферном выведении~\cite{teofilatto}.

Общая модель записывается в инерциальных координатах.
Для объяснения гравитационного разворота и расчётного примера
используется плоская модель невращающейся сферической Земли;
влияние вращения и плоскость орбиты рассмотрены отдельно.
В таблице собраны основные обозначения; остальные вводятся по месту.

\begin{center}
\renewcommand{\arraystretch}{1.12}
\begin{tabularx}{\textwidth}{@{}>{$}l<{$}X@{}}
\toprule
\vect R,\vect V,m & Инерциальные положение, скорость и текущая масса ракеты.\\
r,H,R_E & Радиус $r=\norm{\vect R}$, высота $H=r-R_E$ и радиус Земли.\\
F,c,\vect u & Тяга, скорость истечения $c=g_0I_{\mathrm{sp}}$
и единичное направление тяги.\\
V_r,V_t,\gamma & Радиальная и трансверсальная компоненты скорости;
угол скорости над местным горизонтом.\\
\theta,\beta & Угол тяги над горизонтом и угол
$\beta=\theta-\gamma$ в плоской модели.\\
q_{\mathrm{dyn}},\alpha & Скоростной напор и угол атаки.\\
\vect p,\vect{\mathcal R} & Параметры программы и безразмерная конечная невязка.\\
\bottomrule
\end{tabularx}
\end{center}
$g_0$ --- стандартное ускорение, $I_{\mathrm{sp}}$ --- удельный импульс.
Углы в формулах измеряются в радианах.
Орбита определяется по инерциальным $\vect R,\vect V$.
Тяга и масса рассчитываются на всём активном участке;
$t_f$ обозначает конец всей задачи.

\section{Конечное состояние и целевая орбита}
\label{sec:terminal}

\subsection{Какие условия нужно выполнить}

Высота сама по себе не определяет орбиту.
В центральном поле Земли с гравитационным параметром $\mu$
состояние задаёт удельные энергию
и орбитальный момент:
\begin{equation}
\mathcal E=\frac{V^2}{2}-\frac{\mu}{r},\qquad
\vect h=\vect R\times\vect V,\qquad h=\norm{\vect h}.
\label{eq:invariants}
\end{equation}
Для эллиптической орбиты из них находятся
\begin{equation}
a=-\frac{\mu}{2\mathcal E},\qquad
e=\sqrt{1+\frac{2\mathcal E h^2}{\mu^2}},\qquad
r_p=a(1-e),\quad r_a=a(1+e).
\label{eq:orbital-size}
\end{equation}
Модуль скорости вместе с радиусом определяет энергию,
но момент зависит ещё и от направления скорости.
В плоскости движения
\[
V_r=V\sin\gamma,\qquad V_t=V\cos\gamma,\qquad h=rV_t.
\]
Для круговой орбиты радиуса $r_*$ требуется одновременно
\begin{equation}
\boxed{
r(t_f)=r_*,\qquad V_r(t_f)=0,\qquad
V_t(t_f)=V_*=\sqrt{\frac{\mu}{r_*}}.}
\label{eq:circular-terminal}
\end{equation}
Положение вдоль орбиты здесь свободно.
В пространственной задаче дополнительно задаётся
нужная плоскость; к этому вернёмся в главе о геометрии запуска.

\subsection{Правильная скорость при неправильном направлении}

Возьмём высоту $H_0=650$~км,
$R_E=6378.137$~км и $\mu=398600.4418$~км$^3$/с$^2$.
На радиусе $r_0=7028.137$~км круговая скорость
$V_0=\sqrt{\mu/r_0}=7.530933$~км/с.
Сохраним её модуль, но изменим угол $\gamma_0$.

Энергия остаётся равной $-\mu/(2r_0)$, поэтому $a=r_0$.
Момент равен $\sqrt{\mu r_0}\cos\gamma_0$, откуда
\[
e^2=1-\frac{h^2}{\mu a}
=1-\cos^2\gamma_0=\sin^2\gamma_0.
\]
Следовательно,
\begin{equation}
e=|\sin\gamma_0|,\qquad
H_p=H_0-r_0|\sin\gamma_0|,\qquad
H_a=H_0+r_0|\sin\gamma_0|.
\label{eq:wrong-angle}
\end{equation}
При $\gamma_0=5^\circ$ перигей оказывается на высоте
всего $37.46$~км. Это оскулирующая двухтельная орбита:
в плотной атмосфере реальное движение изменится сопротивлением.
Пример показывает, почему конечное направление скорости
так же существенно, как её модуль.

\coursefigure{geometry-orbits}{Орбиты из одного положения на высоте
650~км при одинаковом модуле скорости и разных углах $\gamma_0$.
Справа --- изменение высоты перигея. Атмосферное движение
здесь не интегрируется.}

\FloatBarrier
\section{Тяга, масса и атмосфера}
\label{sec:dynamics}

Чтобы связать конечную орбиту с возможностями ракеты,
используем знакомую модель конечной тяги и добавим
сопротивление и разделение ступеней.

\subsection{Уравнения движения}

Ракета рассматривается как материальная точка с переменной массой:
\begin{equation}
\begin{aligned}
\dot{\vect R}&=\vect V,\\
\dot{\vect V}&=-\frac{\mu}{r^3}\vect R
+\frac{F}{m}\vect u+\frac{\vect D}{m},\\
\dot m&=-\frac{F}{c},\qquad c=g_0I_{\mathrm{sp}}.
\end{aligned}
\label{eq:motion}
\end{equation}
В нашей расчётной модели $F$ и $c$ постоянны
на каждом активном участке. Масса убывает линейно,
а тяговое ускорение $F/m$ возрастает.
На пассивном участке $F=0$ и расход топлива отсутствует.

Аэродинамические силы зависят от движения относительно воздуха:
\begin{equation}
\vect V_{\mathrm{air}}
=\vect V-\vect\omega_E\times\vect R-\vect w,
\qquad
q_{\mathrm{dyn}}=\frac12\rho(H)V_{\mathrm{air}}^2,
\label{eq:airspeed}
\end{equation}
где $\vect w$ --- ветер относительно вращающейся Земли.
Сопротивление направлено против воздушной скорости:
\begin{equation}
D=q_{\mathrm{dyn}}C_DS,\qquad
\vect D=-\frac12\rho(H)C_DS\,V_{\mathrm{air}}\vect V_{\mathrm{air}}.
\label{eq:drag}
\end{equation}
$C_D$ относится к опорной площади $S$.
Векторная запись определена и при нулевой воздушной скорости:
тогда сопротивление равно нулю.
Подъёмная сила в рассматриваемой модели отсутствует.

Угол атаки $\alpha$ --- угол продольной оси ракеты
к воздушной скорости. Он совпадает с $\beta=\theta-\gamma$
лишь при тяге по оси корпуса и совпадении воздушной
и инерциальной скоростей. Для определения орбиты
всегда используется инерциальная скорость.

\subsection{Топливо и ступени}

Характеристическая скорость определяется интегралом
тягового ускорения на участке от $t_b$ до $t_e$. При постоянном $c$ из
$F\dd t=-c\dd m$ получаем
\begin{equation}
\Delta V_{\mathrm{th}}
=\int_{t_b}^{t_e}\frac{F}{m}\dd t
=c\ln\frac{m_b}{m_e},\qquad
t_e-t_b=\frac{c(m_b-m_e)}{F}.
\label{eq:rocket}
\end{equation}
Первая формула связывает ресурс с массовым отношением,
вторая --- с длительностью при постоянной тяге.
На выведении $\Delta V_{\mathrm{th}}$ не равна
приращению орбитальной скорости: на движение влияют
тяжесть, сопротивление и направление тяги.

В текущую массу входят работающая ступень, верхние ступени
с их топливом и полезная нагрузка. После выгорания первой
ступени её сухая масса ещё остаётся частью ракеты до отделения.
В идеализированной модели разделения
\begin{equation}
\vect R^+=\vect R^-,\qquad
\vect V^+=\vect V^-,\qquad
m^+=m^- -m_{s,1},
\label{eq:separation}
\end{equation}
где $m_{s,1}$ --- сбрасываемая сухая масса.
Положение и скорость непрерывны; меняются масса
и параметры работающего двигателя.

\coursefigure[.68]{geometry-staging}{Массовый состав двухступенчатой ракеты.
Сухая масса первой ступени сбрасывается при разделении;
на активных участках топливо расходуется непрерывно.}

\FloatBarrier
\section{Поворот траектории и программа выведения}
\label{sec:turn}

Ракета начинает движение почти вертикально, а заканчивает
круговое выведение с горизонтальной скоростью.
Чтобы понять этот поворот, запишем уравнения
через модуль скорости и её угол к горизонту.

\subsection{Естественные переменные и уравнения}

Для плоской невращающейся модели положим
\[
\vect V=V(\sin\gamma\,\vect e_r+\cos\gamma\,\vect e_t),
\qquad
\vect u=\sin\theta\,\vect e_r+\cos\theta\,\vect e_t,
\]
где $\vect e_r$ направлен по радиусу наружу,
$\vect e_t$ --- по местному горизонту.
Угол между тягой и скоростью равен $\beta=\theta-\gamma$.

\coursefigure[.60]{geometry-local-frame}{Местный базис, компоненты
скорости и углы $\gamma,\theta,\beta$.
Направление тяги и направление движения могут различаться.}

Кинематика даёт $\dot r=V\sin\gamma$.
Проекция ускорения на скорость содержит
тяговую составляющую $(F/m)\cos\beta$,
сопротивление $-D/m$ и тяжесть $-g\sin\gamma$,
где $g=\mu/r^2$.
Проекция на нормаль в сторону увеличения $\gamma$
описывает поворот скорости в инерциальных осях:
\[
V\frac{\dd(\gamma-\varphi)}{\dd t}
=\frac{F}{m}\sin\beta-g\cos\gamma.
\]
Здесь $\varphi$ --- полярный угол положения,
$\dot\varphi=V\cos\gamma/r$.
Добавляя поворот местного горизонта, получаем
\begin{equation}
\boxed{
\begin{aligned}
\dot r&=V\sin\gamma,\\
\dot\varphi&=\frac{V\cos\gamma}{r},\\
\dot V&=\frac{F}{m}\cos\beta-\frac{D}{m}
                    -\frac{\mu}{r^2}\sin\gamma,\\
\dot\gamma&=\frac{F}{mV}\sin\beta+
 \left(\frac{V}{r}-\frac{\mu}{Vr^2}\right)\cos\gamma,\\
\dot m&=-\frac{F}{c}.
\end{aligned}}
\label{eq:natural}
\end{equation}
Сопротивление не входит в нормальную проекцию,
поскольку направлено вдоль скорости.
При $V=0$ угол не определён, поэтому старт численно
рассчитывается по декартовой системе \eqref{eq:motion}.

\subsection{Начальный и гравитационный разворот}

При строго вертикальном движении
$\gamma=\pi/2$, $\beta=0$ и $\dot\gamma=0$.
В симметричной модели поперечная скорость сама не появится:
нужен начальный разворот тяги.
После набора скорости на вертикальном участке
тяга отклоняется ближе к горизонту,
$\theta<\gamma$, и её нормальная проекция уменьшает $\gamma$.

Затем можно перейти к тяге вдоль скорости.
В нашей невращающейся модели это $\beta=0$;
такой участок называется гравитационным разворотом:
\begin{equation}
\dot\gamma=
\left(\frac{V}{r}-\frac{g}{V}\right)\cos\gamma,
\qquad g=\frac{\mu}{r^2}.
\label{eq:gravity-turn}
\end{equation}
При $V^2<\mu/r$ и $0<\gamma<\pi/2$
угол уменьшается: тяжесть поворачивает траекторию
к местному горизонту.

\coursefigure[.94]{geometry-gravity-turn}{Поворот скорости под действием
тяжести и вращение местного горизонта.
Оба эффекта входят в уравнение для $\gamma$.}

Слагаемое $-g\cos\gamma/V$ описывает поворот скорости
тяжестью, а $V\cos\gamma/r$ --- поворот горизонта.
У поверхности при небольшой скорости вторым можно пренебречь.
Вблизи орбитального режима оно существенно:
при $V^2=\mu/r$ два члена компенсируются,
и горизонтальное движение сохраняется.

\begin{explanation}{Роль начального разворота}
Его момент и величина задают исходные условия последующего
гравитационного разворота. После перехода к тяге вдоль скорости
её направление уже не является независимым управлением.
Выбор начального разворота влияет на высоту, скорость
и возможность дальнейшего довыведения.
\end{explanation}

\subsection{Участки программы}

Типичная рассматриваемая схема включает вертикальный подъём,
начальный разворот, гравитационный разворот,
возможную пассивную дугу и активное довыведение.
В плотной атмосфере малый угол атаки ограничивает боковые нагрузки.
В разреженной атмосфере направление тяги можно выбирать
с большим отклонением от скорости.

\courseplot{geometry-sequence}{Назначение участков выведения.
Схема качественная; высоты и длительности определяются расчётом.}

Пассивная пауза нужна не во всех схемах,
а гравитационный разворот сам по себе не доказывает
оптимальность по топливу или полезной нагрузке.

\FloatBarrier
\section{Потери скорости и ограничения}
\label{sec:losses}

Одна ракета при разных программах получает разные конечные орбиты.
Это удобно объяснять балансом модуля скорости.
Все формулы этой главы относятся к плоской невращающейся
модели без подъёмной силы.

\subsection{Баланс скорости}

В уравнении \eqref{eq:natural} заменим
$(F/m)\cos\beta$ на $F/m-(F/m)(1-\cos\beta)$
и проинтегрируем на всём полёте, включая пассивные участки:
\begin{equation}
V_f-V_0=\Delta V_{\mathrm{th}}
-\Delta V_g-\Delta V_D-\Delta V_{\mathrm{steer}},
\label{eq:loss-balance}
\end{equation}
где
\begin{equation}
\begin{aligned}
\Delta V_{\mathrm{th}}&=\int_{t_0}^{t_f}\frac{F}{m}\dd t,\\
\Delta V_g&=\int_{t_0}^{t_f}\frac{\mu}{r^2}\sin\gamma\,\dd t,\\
\Delta V_D&=\int_{t_0}^{t_f}\frac{D}{m}\dd t,\\
\Delta V_{\mathrm{steer}}&=
 \int_{t_0}^{t_f}\frac{F}{m}(1-\cos\beta)\dd t.
\end{aligned}
\label{eq:loss-integrals}
\end{equation}
Гравитационный интеграл уменьшает скорость на подъёме;
аэродинамический описывает сопротивление;
рулевой --- уменьшение проекции тяги на скорость.
На пассивном подъёме первый из интегралов постоянен,
а гравитационный продолжает расти.
Переход кинетической энергии в потенциальную уже учтён:
добавлять отдельные «затраты на высоту» к этому балансу не нужно.

Крутая траектория быстрее уводит ракету из плотной атмосферы,
но увеличивает гравитационные потери из-за большого $\sin\gamma$.
Ранний поворот уменьшает эту проекцию, но продлевает движение
в плотных слоях. Минимизировать одну потерю недостаточно:
сравниваемые программы должны выполнять конечные условия.

\subsection{Напор и нагрузки}

Скоростной напор $q_{\mathrm{dyn}}=\rho V_{\mathrm{air}}^2/2$
сначала растёт из-за разгона, затем уменьшается
из-за падения плотности.
Кроме предела $q_{\mathrm{dyn}}\le q_{\max}$,
для боковых нагрузок важно произведение
$q_{\mathrm{dyn}}|\alpha|$:
при малом угле нормальная аэродинамическая сила
приближённо пропорциональна ему.
Угол в таком произведении должен быть выражен
в единицах, согласованных с заданным пределом.

Тяговая перегрузка
\begin{equation}
n_T=\frac{F}{mg_0}
\label{eq:thrust-load}
\end{equation}
растёт при постоянной тяге и уменьшающейся массе.
Она характеризует тяговое ускорение;
в атмосфере фактическая продольная нагрузка
включает и сопротивление.
Полная проверка конструкции не сводится
к модели материальной точки.

Таким образом, вместе с орбитой и расходом топлива
нужно проверять максимумы напора и нагрузок,
допустимость управления и остаток топлива.
В расчётном примере этот набор ограничений указан явно.

\FloatBarrier
\section{Пассивный участок и довыведение}
\label{sec:coast}

После работы нижних ступеней ракета может иметь
большую скорость, находясь ещё ниже целевой орбиты.
Пассивная дуга позволяет использовать запас кинетической
энергии для подъёма и выбрать место последнего активного манёвра.

\subsection{Свободный подъём и орбитальная геометрия}

В вакууме без тяги энергия и момент постоянны.
Поэтому сохраняются $a,e,r_p,r_a$,
а текущие высота, скорость и угол траектории меняются.
Из постоянства энергии
\begin{equation}
V_2^2=V_1^2+2\mu\left(\frac1{r_2}-\frac1{r_1}\right).
\label{eq:coast-speed}
\end{equation}
При подъёме скорость уменьшается.
В апогее $V_r=0$, но $V_t$ ещё может быть ниже круговой.
Если апогей ниже цели, одним ожиданием нужной высоты
достичь нельзя.

Для промежуточной эллиптической орбиты с
$r_a=r_*$ и $a_{\mathrm{tr}}=(r_p+r_*)/2$
импульсная оценка довыведения в апогее равна
\begin{equation}
V_a^-=\sqrt{\mu\left(\frac2{r_*}
-\frac1{a_{\mathrm{tr}}}\right)},\qquad
\Delta V_{\mathrm{circ}}=\sqrt{\frac{\mu}{r_*}}-V_a^-.
\label{eq:coast-circularize}
\end{equation}
Трансверсальный импульс поднимает перигей до $r_*$.

\coursefigure{geometry-coast-insertion}{Пассивная эллиптическая дуга
и идеальное импульсное довыведение в апогее.}

\subsection{Выбор начала активного довыведения}

При конечной тяге положение и скорость меняются во время
самого довыведения. Поэтому начало манёвра в точке
идеального импульса в общем случае не создаёт круговую орбиту.
Начало выбирается так, чтобы одновременно выполнить
три условия \eqref{eq:circular-terminal}.

Например, перед последним манёвром $V_r$ может быть положительной.
За время набора $V_t$ тяжесть и радиальная проекция тяги
должны довести $V_r$ до нуля.
Длительность предшествующей паузы определяет,
с какими $r,V_r,V_t$ начнётся этот процесс.
Именно поэтому её подбирают вместе с длительностью
и направлением последнего активного участка.

\FloatBarrier
\section{Подбор программы выведения}
\label{sec:targeting}

После задания программы выполняется прогноз
и получается некоторая конечная орбита.
Параметры программы $\vect p$ уточняются
методом пристрелки из предыдущей лекции.
Для круговой цели используем одну безразмерную невязку:
\begin{equation}
\vect{\mathcal R}(\vect p)=
\begin{pmatrix}
(r_f-r_*)/r_*\\
V_{r,f}/V_*\\
(V_{t,f}-V_*)/V_*
\end{pmatrix},
\qquad
\norm{\vect{\mathcal R}}_2\le\varepsilon.
\label{eq:residual}
\end{equation}
Масштабы $r_*$ и $V_*=\sqrt{\mu/r_*}$ приводят
компоненты к согласованной форме.
Если конечное время свободно, длительность манёвра
может входить в $\vect p$.

Для текущей программы матрица
$\partial\vect{\mathcal R}/\partial\vect p$
вычисляется центральными разностями с повторными прогнозами.
Линеаризация даёт ньютоновскую поправку:
\begin{equation}
\frac{\partial\vect{\mathcal R}}{\partial\vect p}
\,\delta\vect p=-\vect{\mathcal R}(\vect p),\qquad
\vect p_{\mathrm{new}}=\vect p+\eta\,\delta\vect p,
\quad 0<\eta\le1.
\label{eq:newton}
\end{equation}
После обновления снова рассчитывается вся траектория.
Если невязка растёт либо параметры нарушают ограничения,
шаг уменьшается. Итерации продолжаются до выполнения
допуска в \eqref{eq:residual}; отдельно проверяются
нагрузки и запас топлива.

\textbf{Пристрелка обеспечивает конечные условия.}
Максимум остатка топлива или полезной нагрузки требует
отдельной оптимизации. В примере ищется допустимая программа.

\FloatBarrier
\section{Вращение Земли и геометрия запуска}
\label{sec:launch-geometry}

В пространственной задаче плоскость орбиты согласуют с местом,
направлением и временем старта.

\subsection{Вращение Земли и две скорости}

Ракета на площадке и атмосфера вращаются вместе с Землёй.
Поэтому начальная инерциальная скорость ракеты и её скорость
относительно воздуха без ветра равны
\begin{equation}
\begin{aligned}
 \vect V_0&=\vect\omega_E\times\vect R_0,
 \qquad V_{E,0}=\omega_E R_E\cos\phi,\\
 \vect V_{\mathrm{air}}&=\vect V-\vect\omega_E\times\vect R.
\end{aligned}
\label{eq:launch-initial-velocity}
\end{equation}
Здесь $\phi$ --- геоцентрическая широта. Скорость площадки направлена
на восток и на экваторе составляет около $465$~м/с.
Орбиту определяет $\vect V$, аэродинамику --- $\vect V_{\mathrm{air}}$.
Из векторного сложения видно, что различаются их модули и азимуты.

\courseplot[.55]{geometry-rotation}{Инерциальная и воздушная скорости
без ветра. Скорость атмосферы направлена на восток;
$\vect V=\vect V_{\mathrm{air}}+\vect\omega_E\times\vect R$.}

\subsection{Азимут и наклонение}

Азимут $A$ отсчитываем от севера к востоку: $A=0$ соответствует
движению на север, $A=90^\circ$ --- на восток.
Пусть $\vect e_E$ и $\vect e_N$ --- местные направления востока и севера,
а $\vect k$ --- единичный вектор земной оси.
Единичное горизонтальное направление движения обозначено
на рисунке через $\vect e_A$.
Радиальная скорость не вносит вклада в орбитальный момент,
поэтому нормаль орбиты определяется горизонтальным направлением движения:
\begin{equation}
 \vect n=\frac{\vect R\times\vect V}{\norm{\vect R\times\vect V}}
        =\sin A\,\vect e_N-\cos A\,\vect e_E.
 \label{eq:orbit-normal-azimuth}
\end{equation}
В местном базисе $\vect k=\sin\phi\,\vect e_r+\cos\phi\,\vect e_N$.
Скалярное произведение нормали на земную ось даёт
\begin{equation}
 \boxed{\cos i=\vect n\cdot\vect k=\cos\phi\sin A.}
 \label{eq:launch-azimuth-inclination}
\end{equation}
Так широта и направление движения задают наклонение орбиты.
Здесь $A$ --- \emph{инерциальный} азимут, а $\phi$ относится
к точке, в которой определены положение и скорость.
Подстановка широты космодрома даёт исходную геометрическую оценку
для запуска в плоскость, проходящую через него~\cite{teofilatto}.

\courseplot[.96]{geometry-launch-plane}{Связь местного направления
движения с плоскостью орбиты: проекция её нормали на земную ось
равна $\cos i=\cos\phi\sin A$.}

\subsection{Время старта и плоскость орбиты}

Одно наклонение ещё не фиксирует плоскость: при том же $i$
возможны разные долготы восходящего узла $\Omega$.
Космодром вращается вместе с Землёй, а целевая плоскость
фиксирована в инерциальных осях. Поэтому для прямого запуска
нужно выбрать момент, когда космодром окажется в этой плоскости:
\begin{equation}
 \boxed{\vect n_*\cdot\vect R_s(t_0)=0.}
 \label{eq:launch-window-plane}
\end{equation}
Здесь $\vect n_*$ --- нормаль целевой плоскости,
$\vect R_s(t)$ --- инерциальное положение космодрома.
На рисунке этому условию соответствуют пересечения суточной
параллели с плоскостью орбиты; $\Lambda(t)$ обозначает
инерциальную долготу космодрома.

\courseplot[.92]{geometry-launch-window}{Геометрия выбора времени
старта. В экваториальной проекции космодром движется по окружности;
положения $A$ и $B$ соответствуют его прохождению через целевую плоскость.}

\FloatBarrier
\begin{explanation}{Что задаёт геометрия запуска}
Широта и инерциальный азимут определяют наклонение,
а время старта позволяет выбрать ориентацию плоскости.
Это исходная геометрическая оценка: полная программа выведения
должна обеспечить конечное попадание в плоскость с учётом
вращения Земли и движения ракеты во время полёта.
\end{explanation}
\FloatBarrier


\FloatBarrier
\section{Расчёт выведения учебной двухступенчатой ракеты}
\label{sec:numerical-ascent}
% Generated by ascent_numerics.py; units documented in source.
\newcommand{\AscentTOne}{134{,}491}
\newcommand{\AscentRaise}{147{,}286}
\newcommand{\AscentCoast}{301{,}229}
\newcommand{\AscentCircularize}{83{,}044}
\newcommand{\AscentFinalTime}{666{,}050}
\newcommand{\AscentFinalMass}{12{,}245}
\newcommand{\AscentReserve}{3{,}245}
\newcommand{\AscentCharSpeed}{10{,}048}
\newcommand{\AscentGravityLoss}{2303{,}188}
\newcommand{\AscentDragLoss}{36{,}440}
\newcommand{\AscentSteeringLoss}{177{,}407}
\newcommand{\AscentMaxQ}{31{,}508}
\newcommand{\AscentMaxQAlpha}{0{,}723}
\newcommand{\AscentMaxLoad}{6{,}662}
\newcommand{\AscentPositionCheck}{1.20\cdot10^{-7}}
\newcommand{\AscentVelocityCheck}{3.32\cdot10^{-10}}
\newcommand{\AscentEccentricity}{8.28\cdot10^{-14}}
\newcommand{\AscentResidual}{4.66\cdot10^{-14}}
\newcommand{\AscentBalanceError}{-6.22\cdot10^{-12}}


\subsection{Исходные данные и программа}

Выведем нагрузку массой $5$~т на круговую орбиту высотой $650$~км.
Цель примера --- проследить, как программа тяги формирует конечную
орбиту. Это учебная ракета; расчёт выполняется в плоскости вокруг
невращающейся сферической Земли. Примем
\[
 \mu=398600{,}4418\ \text{км}^3/\text{с}^2,\quad
 R_E=6378{,}137\ \text{км},\quad g_0=9{,}80665\ \text{м}/\text{с}^2.
\]
На старте $\vect R_0=(R_E,0)^\top$, $\vect V_0=0$;
целевые значения $r_*=7028{,}137$~км и $V_*=7{,}5309329$~км/с.

\begin{table}[htbp]
\centering
\caption{Характеристики учебной ракеты.}
\label{tab:ascent-vehicle}
\renewcommand{\arraystretch}{1.12}
\begin{tabularx}{\textwidth}{@{}Xrr@{}}
\toprule
Параметр & Первая ступень & Верхняя ступень\\\midrule
Масса топлива, т & $160$ & $45$\\
Сухая масса, т & $16$ & $4$\\
Тяга, МН & $3{,}5$ & $0{,}8$\\
Удельный импульс, с & $300$ & $450$\\
\bottomrule
\end{tabularx}
\end{table}

Начальная масса равна $230$~т. Первая ступень расходует всё топливо
за $\AscentTOne$~с: к разделению остаётся $70$~т, после сброса
её сухой массы --- $54$~т. Верхний двигатель допускает выключение
и повторный запуск. Тяга и удельный импульс каждой ступени постоянны;
переходные процессы двигателя и поворота корпуса не моделируются.

Атмосферу зададим экспонентой
\begin{equation}
 \rho(H)=1{,}225\exp\left(-\frac{H}{7{,}2\ \text{км}}\right)
 \ \text{кг}/\text{м}^3,\qquad C_D=0{,}3,\quad S=10\ \text{м}^2.
 \label{eq:numerical-atmosphere}
\end{equation}
В сопротивлении скорость переводится в м/с.
Атмосфера неподвижна, поэтому воздушная и инерциальная скорости
совпадают; подъёмная сила не учитывается.

Программа состоит из шести участков:
\begin{enumerate}
\item Вертикальная тяга в течение $10$~с.
\item Линейное уменьшение $\theta$ от $90^\circ$ до $70^\circ$
за следующие $25$~с --- начальный разворот.
\item Тяга вдоль скорости, $\theta=\gamma$, до выгорания первой ступени.
\item После разделения --- первый активный участок верхней ступени
при $\theta=45^\circ$, формирующий переходную дугу.
\item Пассивный подъём при выключенном двигателе.
\item Трансверсальное довыведение, $\theta=0$, до получения
круговой скорости при $V_r=0$.
\end{enumerate}
Изменения направления на границах идеализируются мгновенными;
положение и скорость остаются непрерывными.

\courseplot[.98]{ascent-trajectory}{Траектория с выделением участков.
Дальность равна $R_E\varphi$, где $\varphi$ --- угол от радиуса точки
старта. Масштабы высоты и дальности различны.}

\subsection{Пристрелка длительностей}

Изменяем длительности первого верхнего активного участка,
пассивной паузы и довыведения:
\[
 \vect p=(\tau_{\mathrm a},\tau_{\mathrm c},\tau_{\mathrm d})^\top.
\]
Доступное топливо ограничивает сумму активных участков:
\[
 \tau_{\mathrm a}+\tau_{\mathrm d}\le248{,}231\ \text{с}.
\]
Используем невязку~\eqref{eq:residual} с допуском
$\norm{\vect{\mathcal R}}_2<10^{-11}$.
От начального приближения
$\vect p^{(0)}=(150,300,80)^\top$~с выполнены три ньютоновские
поправки с производными по центральным разностям.
Каждая оценка невязки включает полный прогноз от старта.
Получено
\begin{equation}
 \boxed{
 \tau_{\mathrm a}=\AscentRaise\ \text{с},\qquad
 \tau_{\mathrm c}=\AscentCoast\ \text{с},\qquad
 \tau_{\mathrm d}=\AscentCircularize\ \text{с}.}
 \label{eq:numerical-solution}
\end{equation}
Конечный момент $t_f=\AscentFinalTime$~с.
Здесь параметры округлены; для расчёта используются
их полные значения из сохранённых данных.

\courseplot[.88]{ascent-convergence}{Уменьшение нормы безразмерной
невязки за три поправки. Номер $0$ соответствует исходной программе.}

\subsection{Формирование орбиты}

После первого активного участка верхней ступени высота составляет
около $283$~км, а радиальная скорость --- $2{,}10$~км/с.
Далее ракета поднимается без тяги за счёт уменьшения
кинетической энергии.

\courseplot[.98]{ascent-state-history}{Высота, компоненты скорости
и углы на общей временной шкале. На пассивном участке линия $\theta$
не показана. При $V=0$ угол $\gamma$ не определён; начальное значение
$90^\circ$ обозначает вертикальный предел.}

Довыведение начинается до апогея: $H=640{,}845$~км,
$V_r=0{,}2973$~км/с и $V_t=3{,}9991$~км/с.
При трансверсальной тяге и пренебрежимо малом сопротивлении
\begin{equation}
 \dot V_r=\frac{V_t^2}{r}-\frac{\mu}{r^2}<0
 \quad\text{при}\quad V_t<\sqrt{\frac{\mu}{r}}.
 \label{eq:numerical-radial-deceleration}
\end{equation}
Поэтому положительная радиальная скорость уменьшается без
направленной вниз тяги. Пауза подобрана так, чтобы $V_r$
достигла нуля одновременно с получением круговой скорости.

\courseplot[.94]{ascent-final-insertion}{Довыведение: радиальная
скорость и дефицит поперечной скорости одновременно стремятся к нулю.}

Перед последним манёвром апогей переходной орбиты равен
$648{,}455$~км, а перигей лежит внутри Земли.
Отрицательная высота перигея описывает продолжение оскулирующей
коники; оставлять аппарат на такой траектории нельзя.
На пассивной дуге апсиды практически постоянны.
Довыведение поднимает и апогей, и перигей до $650$~км.
Таким образом, последний манёвр не сводится к импульсу точно
в апогее целевой высоты.

\courseplot[.98]{ascent-osculating-orbit}{Изменение апсид и удельной
механической энергии. Постоянный участок соответствует пассивной дуге;
в конце апогей и перигей совпадают.}

\FloatBarrier
\subsection{Топливо, нагрузки и потери}

Конечная масса равна $\AscentFinalMass$~т, включая $9$~т сухой
верхней ступени и нагрузки. Остаток топлива --- $\AscentReserve$~т.
На активных участках масса убывает непрерывно; разделение создаёт
её скачок. Уменьшение массы при постоянной тяге увеличивает $n_T$.

\courseplot[.98]{ascent-propulsion}{Масса, тяга и тяговая перегрузка.
Сброс массы при разделении не создаёт скачка скорости.}

После пристрелки проверены иллюстративные пределы
$q_{\mathrm{dyn}}\le50$~кПа и $n_T\le8$.
Максимумы на сетке с интервалом не более $1$~с составляют
$\AscentMaxQ$~кПа и $\AscentMaxLoad$ соответственно.
При тяге по продольной оси в начальном развороте
$|\alpha|_{\max}\simeq2{,}68^\circ$,
а $(q_{\mathrm{dyn}}|\alpha|)_{\max}=\AscentMaxQAlpha$~кПа$\cdot$рад.
На верхнем участке рассогласование тяги и скорости достигает
$20{,}4^\circ$, но напор уже не превышает $0{,}323$~кПа.
Это объясняет, почему одинаковое отклонение направления
создаёт разные аэродинамические нагрузки на разных высотах.
Отдельный расчёт прочности здесь не выполняется.

\courseplot[.98]{ascent-aeroloads}{Напор и произведение
$q_{\mathrm{dyn}}|\alpha|$ на первой ступени; угол выражен в радианах.
После перехода к тяге вдоль скорости $\alpha=0$.}

Скоростной баланс даёт
\begin{equation}
 \underbrace{10{,}047968}_{\Delta V_{\mathrm{th}}}
 -\underbrace{2{,}303188}_{\Delta V_g}
 -\underbrace{0{,}036440}_{\Delta V_D}
 -\underbrace{0{,}177407}_{\Delta V_{\mathrm{steer}}}
 =7{,}530933\ \text{км}/\text{с}.
 \label{eq:numerical-speed-budget}
\end{equation}
Гравитационный интеграл растёт и на пассивной восходящей дуге:
характеристическая скорость постоянна, а текущая скорость уменьшается.
Сопротивление в основном влияет на первую ступень;
рулевые потери накапливаются прежде всего при $\theta=45^\circ$.

\courseplot[.98]{ascent-losses}{Накопленные потери и различие
характеристической и текущей скоростей. Подъём потенциальной энергии
уже учтён в балансе через гравитационный интеграл.}

\FloatBarrier
\subsection{Влияние начального разворота}

Сохраним найденные длительности и изменим только начальный разворот:
$16^\circ$ или $24^\circ$ вместо $20^\circ$.
Расход топлива и характеристическая скорость останутся одинаковыми.

Крутая программа заканчивается на высоте $793$~км,
но её перигей около $-126$~км: большая высота не обеспечила
орбитальный полёт. Пологая программа заканчивается на высоте
$501$~км с отрицательной $V_r$ и получает эллиптическую орбиту
с апсидами $447$ и $1326$~км. Обе программы не выполняют
круговую цель.

\courseplot[.98]{ascent-program-comparison}{Разные результаты трёх
начальных разворотов при одинаковых последующих длительностях
и расходе топлива. Соседние программы повторно не пристреливались.}

\begin{explanation}{Результат примера}
Пристрелка обеспечила круговую орбиту $650$~км с положительным
остатком топлива и выполнением выбранных проверок напора и $n_T$.
Она не минимизировала расход и не максимизировала полезную нагрузку.
Сравнение программ показывает роль управления при одинаковом
запасе характеристической скорости.
\end{explanation}

Результат подтверждён повторным интегрированием с более строгими
допусками и проверками массового и скоростного балансов.
Расчёт и рисунки воспроизводятся скриптом
\texttt{10/images/ascent\_numerics.py}; параметры, траектории и
проверки сохранены рядом в JSON, CSV и NPZ.


\FloatBarrier

\begin{thebibliography}{9}
\addcontentsline{toc}{section}{Литература}
\bibitem{teofilatto}
Teofilatto P. \textit{Trajectories of Launch and Reentry Systems}.
Springer Aerospace Technology. Springer, 2025.
\href{https://doi.org/10.1007/978-3-031-75760-0}
{DOI: 10.1007/978-3-031-75760-0}.
Разд.~1.2, 1.8--1.9, 2.2--2.7, 3.2--3.4:
ступени, аэродинамика, геометрия запуска,
атмосферное движение, потери, пассивные участки и довыведение.
\bibitem{nasa-rocket}
NASA Glenn Research Center.
\href{https://www1.grc.nasa.gov/beginners-guide-to-aeronautics/ideal-rocket-equation/}{\textit{Ideal Rocket Equation}};
\href{https://www1.grc.nasa.gov/beginners-guide-to-aeronautics/dynamic-pressure-2/}{\textit{Dynamic Pressure}}.
Определения характеристической скорости и скоростного напора.
\end{thebibliography}
\end{document}
